% Part: second-order-logic
% Chapter: sol-and-set-theory
% Section: introduction

\documentclass[../../../include/open-logic-section]{subfiles}

\begin{document}

\olfileid{sol}{set}{int}

\olsection{مقدمه}

از آنجا که منطق مرتبهٔ دوم می‌تواند افزون بر توابع بر زیرمجموعه‌های
دامنه نیز کمیت‌گذاری کند، انتظار می‌رود دست‌کم بخشی از نظریهٔ مجموعه‌ها
را بتوان در منطق مرتبهٔ دوم صورت‌بندی کرد. مقصودمان از «صورت‌بندی» این
است که می‌توان ویژگی‌ها و گزاره‌های مجموعه‌شناختی را در منطق مرتبهٔ دوم
بیان کرد، و این کار بدون هیچ واژگان ویژه و غیرمنطقی برای مجموعه‌ها
(برای نمونه، !!{predicate} عضویت در نظریهٔ مجموعه‌ها) ممکن است. برای
نمونه، می‌توانیم اجتماع و اشتراک مجموعه‌ها و رابطهٔ زیرمجموعه‌بودن را
تعریف کنیم، اما همچنین می‌توانیم اندازهٔ مجموعه‌ها را مقایسه کنیم و
نتایجی چون قضیهٔ کانتور را بیان کنیم.

\end{document}
